% Original: PDF strana 8, gornji slajd; prvi deo odobrene podele.
\slajdid{02-s015}
\begin{frame}[label=02-s015]{Uslov za red sa logičkom nulom}
\begin{columns}[T,onlytextwidth]
\begin{column}{.66\textwidth}
% element: 02-s015:e01
Funkciju možemo formirati polazeći od redova u kojima je \(F=0\).
\par\vspace{3mm}
% element: 02-s015:e02
\Oznaka{Prvi red: \((C,B,A)=(0,0,0)\)}\vspace{2mm}
Uslov je ispunjen kada je \(C=0\) \alert{I} \(B=0\) \alert{I} \(A=0\).
\par\vspace{2mm}
Proizvod \(\bar C\bar B\bar A\) tada je 1. Za činilac koji treba da bude \alert{0} uzimamo njegovu negaciju:
\[\textcolor{DEcrvena}{\overline{\bar C\bar B\bar A}=C+B+A}.\]
\Rezultat{NI uslov sa aktivnom nulom odgovara ILI funkciji: zbir je 0 samo kada su svi njegovi ulazi 0.}
\end{column}
\begin{column}{.29\textwidth}\centering
\Oznaka{Funkcionalna tabela}\vspace{3mm}
{\fontsize{9}{11}\selectfont\begingroup
\setlength{\tabcolsep}{2pt}
\renewcommand{\arraystretch}{1.04}
\par\noindent
\begin{tabularx}{\linewidth}{*{4}{>{\centering\arraybackslash}X}}
\toprule C&B&A&F\\\midrule
\textcolor{DEcrvena}{0}&\textcolor{DEcrvena}{0}&\textcolor{DEcrvena}{0}&\textcolor{DEcrvena}{0}\\
\textcolor{DEcrvena}{0}&\textcolor{DEcrvena}{0}&\textcolor{DEcrvena}{1}&\textcolor{DEcrvena}{0}\\
\textcolor{DEcrvena}{0}&\textcolor{DEcrvena}{1}&\textcolor{DEcrvena}{0}&\textcolor{DEcrvena}{0}\\
0&1&1&1\\
\textcolor{DEcrvena}{1}&\textcolor{DEcrvena}{0}&\textcolor{DEcrvena}{0}&\textcolor{DEcrvena}{0}\\
1&0&1&1\\
1&1&0&1\\
\textcolor{DEcrvena}{1}&\textcolor{DEcrvena}{1}&\textcolor{DEcrvena}{1}&\textcolor{DEcrvena}{0}\\\bottomrule
\end{tabularx}\par
\endgroup
}
\end{column}
\end{columns}
\end{frame}
